% Unofficial Yale Poster template.
% A fork of the UMich template https://www.overleaf.com/latex/templates/university-of-michigan-umich-poster-template/xpnqzzxwbjzc
% which is fork of the MSU template https://www.overleaf.com/latex/templates/an-unofficial-poster-template-for-michigan-state-university/wnymbgpxnnwd
% which is a fork of https://www.overleaf.com/latex/templates/an-unofficial-poster-template-for-new-york-university/krgqtqmzdqhg
% which is a fork of https://github.com/anishathalye/gemini
% also refer to https://github.com/k4rtik/uchicago-poster



\documentclass[final]{beamer}

% ====================
% Packages
% ====================

\usepackage[T1]{fontenc}
 \usepackage[utf8]{luainputenc}
\usepackage{lmodern}
\usepackage[size=custom, width=122,height=91, scale=1.2]{beamerposter}
\usetheme{gemini}
\usecolortheme{msu}
\usepackage{graphicx}
\usepackage{booktabs}
\usepackage{tikz}
\usepackage{pgfplots}
\pgfplotsset{compat=1.14}
\usepackage{anyfontsize}
\usepackage{wrapfig}
\usepackage{amsmath}
\usepackage[numbers]{natbib}

% ====================
% Lengths
% ====================

% If you have N columns, choose \sepwidth and \colwidth such that
% (N+1)*\sepwidth + N*\colwidth = \paperwidth
\newlength{\sepwidth}
\newlength{\colwidth}
\setlength{\sepwidth}{0.025\paperwidth}
\setlength{\colwidth}{0.3\paperwidth}

\newcommand{\separatorcolumn}{\begin{column}{\sepwidth}\end{column}}
\renewcommand{\bibfont}{\small} 

% ====================
% Title
% ====================

\title{Low-dimensional coexistence in complex microbial ecosystems}

\author{Karthik Srinivasan\textsuperscript{1}, Germ\'an Plata\textsuperscript{2}, Purushottam Dixit\textsuperscript{1}}

\institute[shortinst]{\textsuperscript{1}Department of Biomedical Engineering, Yale University, \textsuperscript{2}Computational Sciences, BiomEdit, LLC.}

% ====================
% Footer (optional)
% ====================

\footercontent{
  % \href{https://github.com}{Github: https://github.com} 
  \hfill
Poster Session \hfill
  \href{mailto:youremail@yale.edu}{karthik.srinivasan@yale.edu}}
% (can be left out to remove footer)

% ====================
% Logo (optional)
% ====================

% use this to include logos on the left and/or right side of the header:
% Left: institution
 \logoleft{\includegraphics[height=8cm]{logos/logo.png}}
% Right: funding agencies and other affilations 
\logoright{\includegraphics[height=5.5cm]{figures/attribution and qr.png}}
% ====================
% Body
% ====================

\begin{document}

\begin{frame}[t]
\begin{columns}[t]
\separatorcolumn

\begin{column}{\colwidth}
    
        
    \begin{block}{Niche theory predicts niche dimensionality to drive diversity}
        
        \begin{minipage}{0.6\textwidth}
            % \vspace{20pt}
            \begin{itemize}
                \item Current physical models treat microbiomes as disordered systems with random, high-dimensional interactions (akin to spin glasses).
                \item This core assumption of high-dimensional coexistence remains untested empirically.
                \item Niche theory predicts higher niche dimensionality, $\eta_D$, drives diversity by reducing niche overlap, but estimating $\eta_D$ is difficult.
                \item We take an effective approach to estimate $\eta_D$ directly from abundance profile data.
            \end{itemize}
        \end{minipage}
        \hfill
        \begin{minipage}{0.4\textwidth}
            \begin{figure}
                \centering
                \hspace{-20pt} 
                \includegraphics[width=\linewidth]{figures/Niche overlap blobs.jpg}
            \end{figure}
            \scriptsize{Pianka, Eric R. Evolutionary ecology. Eric R. Pianka, 2011.}
        \end{minipage}
        
    \end{block}
    \vspace{-20pt}
    \begin{block}{Latent variable JSDMs can be fit to data and is interpretable}
        \begin{minipage}{0.5\textwidth}
        \vspace{20pt}
        \begin{itemize}
            \item JSDMs relate species distribution to environmental factors and estimate $\eta_D$.
            \vspace{10pt}
            \item Latent variable formulization of JSDMs decompose microbiomes into interpretable sample specific and species specific embeddings.
            \vspace{10pt}
            \item We define $\eta_D$ as the characteristic environmental dimensionality that explains composition.
        \end{itemize}
        \end{minipage}
        \hfill
        \begin{minipage}{0.5\textwidth}
            \begin{figure}
                \centering
                \hspace{-20pt}
                \includegraphics[width=\linewidth]{figures/model as matrix factorization.png}
            \end{figure}
        \end{minipage}
        \begin{figure}
            \centering
            \vspace{20pt}
            \includegraphics[width=\linewidth]{figures/model interpretability.png}
        \end{figure}
        \vspace{-10pt}
        
    \end{block}


    \begin{block}{Estimated dimensionality reflects environmental complexity and follows niche dimensionality hypothesis}
         \begin{itemize}
            \item Inferred $\eta_D$ correlates with the dimensionality of metabolic environment, validating our abundance only approach.
            \vspace{10pt}
            \item Consistent with the macroecological niche dimensionality hypothesis, $\eta_D$ correlates with niche overlap and biodiversity.
            \vspace{10pt}
        \end{itemize}
         \begin{figure}
            \centering
            \includegraphics[width=1\linewidth]{figures/Paper_fig_2.png}
        \end{figure}
       
     
  \end{block}
  
\end{column}
\separatorcolumn
\begin{column}{\colwidth}

  \begin{block}{Consumer resource models are mechanistic and widely used}
  \vspace{20pt}
    \begin{minipage}{0.5\linewidth}
        \begin{itemize}
            \item The consumer resource framework uses environmental mediation of bacterial interaction to model microbiomes.
            \vspace{10pt}
            \item We incorporate two types of dimensional constraints that can theoretically allow species to exist in low-dimensional niches.
        \end{itemize}
    \end{minipage}
    \hfill
    \begin{minipage}{0.5\linewidth}
        \begin{align*}
                &\frac{1}{n_o}\frac{dn_o}{dt} = \sum_{k=1}^{K} r_k(t)\theta_{ko} - \delta_o \\
                &\frac{dr_k}{dt} = J_k(t) - \sum_{o=1}^O r_k(t) \frac{\theta_{ko}}{w_k} n_o(t)
            \end{align*}
    \end{minipage}    
    \begin{figure}
        \centering
        % \hspace{-150pt}
        \vspace{10pt}
        \includegraphics[clip, width=0.9\textwidth]{figures/Model.png}
        \vspace{10pt}
    \end{figure}
     \begin{itemize}
        \item \textbf{Intrinsic tradeoffs:} Modeled as simplex constraints on species preference vectors, placing them on equal footing.
        \item \textbf{Extrinsic tradeoffs:} Modeled as low-rank correlations in nutrient inflow vectors that reduce the effective degrees of freedom in the environment.
     \end{itemize}
  \end{block}

  \begin{block}{Tradeoffs explains how $\eta_D$ relates with community structure}
        \begin{itemize}
            \vspace{10pt}
            \item Tradeoffs successfully reproduce the empirical relationships between community structure, niche overlap, and niche dimensionality.
        \end{itemize}
        \begin{figure}
        \centering
            \centering
            \vspace{10pt}
            \includegraphics[width=\linewidth]{figures/in silico results.png}
            \vspace{-20pt}
        \end{figure}
        \begin{itemize}
            \item \textbf{Crucially, both intrinsic (simplex constraints) and extrinsic (low-rank correlations) tradeoffs are strictly required} to explain the full set of empirical trends.
        \end{itemize}
    \end{block}

  

\end{column}
\separatorcolumn
\begin{column}{\colwidth}.

    \begin{block}{Tradeoffs explain low $\eta_D$ of microbiomes}

    \begin{itemize}
        \item Across $\sim 200$ real microbiome datasets obtained from the Microbiomap project we see that microbiomes are highly low dimensional.
        \item \textbf{Disordered consumer-resource models tuned to match empirical diversity predict median $\eta_D \approx 13$, directly contradicting the empirically observed median $\eta_D \approx 4$.}
        \item In simulations, extrinsic tradeoffs (correlated resource supplies) impose stronger dimensional constraints on the effective niche space than intrinsic metabolic tradeoffs alone.
    \end{itemize}
    \begin{figure}
        \centering
        \includegraphics[width=\linewidth]{figures/BoxPlot Niche dimensionality.png}
    \end{figure}

  \end{block}
    


  \begin{block}{Implications for Statistical Physics Models}
    \begin{itemize}
         \item Our empirical estimates directly challenge theoretical models that draw interactions or consumer preferences i.i.d. from fixed distributions.
         \vspace{5pt}
         \item Coexistence in complex microbial ecosystems is not governed by effectively infinite, independent ecological axes.
         \vspace{5pt}
         \item To accurately capture macroscopic ecological dynamics, future statistical physics models must explicitly incorporate dimensional constraints.
         \vspace{5pt}
         \item \textbf{Key Takeaway:} Incorporating low-rank interaction matrices (in Generalized Lotka-Volterra systems) and correlated resource structures (in CRMs) is essential to model realistic microbiome phase spaces.
     \end{itemize}
    
  \end{block}
    \vspace{-10pt}
  \begin{block}{References}
  % \vspace{-5pt}
    \nocite{*}
    \footnotesize{\bibliographystyle{unsrtnat}\bibliography{poster}}

  \end{block}

\end{column}

\separatorcolumn
\end{columns}
\end{frame}

\end{document}
